\section*{Capítulo \ref{ch:b3:virology} --- Virologia}

\begin{solution}{exo:b3:virology:1}
Poliovírus: IV, RNA de fita positiva --- uma RNA-polimerase dependente
de RNA. Influenza: V, RNA de fita negativa --- sua própria
RNA-polimerase, carregada no \omterm{def:b3:virology:virus}{vírion}. HIV: VI --- transcriptase reversa
(e integrase). Herpes simples: I, DNA de fita dupla --- o hospedeiro
poderia em princípio fornecer tudo, mas ele traz sua própria
DNA-polimerase e uma timidina-cinase. Hepatite B: VII --- uma
transcriptase reversa para copiar seu RNA pré-genômico em DNA.
Rotavírus: III, RNA de fita dupla --- uma RNA-polimerase dependente de
RNA dentro da partícula.
\end{solution}

\begin{solution}{exo:b3:virology:2}
O DNA do fago ($^{32}$P) entrou nas bactérias e a proteína ($^{35}$S)
ficou do lado de fora e pôde ser arrancada; a prole do fago herdou o
$^{32}$P. Logo, o DNA é o que o fago injeta e o que dirige a produção
de fagos novos. Se a proteína carregasse as instruções, o $^{35}$S
teria sido encontrado dentro e passado à prole.
\end{solution}

\begin{solution}{exo:b3:virology:3}
$84/(0.1\times 10^{-7}) = 8.4\times 10^{9}$ unidades formadoras de
placa por mL. O microscópio conta cada partícula, inclusive \omterm{def:b3:virology:virus}{capsídeos}
vazios, partículas de genoma defeituoso e as que não conseguem aderir
ou injetar; só as partículas que completam uma infecção formam placas.
\end{solution}

\begin{solution}{exo:b3:virology:4}
Adesão: inibidores de entrada (maraviroque), anticorpos
neutralizantes. Entrada/fusão: inibidores de fusão, bloqueadores da
acidificação endossomal. Desnudamento: fármacos que se ligam ao
\omterm{def:b3:virology:virus}{capsídeo} (pleconaril; lenacapavir). Expressão e replicação: \omterm{def:b3:virology:drugs}{análogos de
nucleosídeo}, inibidores de transcriptase reversa e de integrase,
interferon, \omterm{def:b3:rna-regulation:rnai}{interferência por RNA}. Montagem: \omterm{def:b3:virology:drugs}{inibidores de protease}
(poliproteínas não clivadas não se montam). Liberação: inibidores de
neuraminidase (influenza), e linfócitos~T citotóxicos que matam a
célula antes da liberação.
\end{solution}

\begin{solution}{exo:b3:virology:5}
Uma proteína de \qty{30}{kDa} tem $30\,000/110 = 273$ resíduos, $818$
nucleotídeos: \qty{0.82}{kb}, \qty{18}{\%} de um genoma de
\qty{4.5}{kb}. Codificar as $180$ cópias separadamente exigiria
\qty{147}{kb}, trinta e três vezes o genoma inteiro; a própria proteína
do \omterm{def:b3:virology:virus}{capsídeo} (\qty{5.4}{MDa}) pesa quatro vezes mais que o genoma
(\qty{1.5}{MDa}).
\end{solution}

\begin{solution}{exo:b3:virology:6}
$V = 10^{5}(3e^{-0.5t} - 0.5e^{-3t})/2.5$: $t = 0.5$: $8.9\times
10^{4}$; $t = 2$: $4.4\times 10^{4}$; $t = 10$: $8\times 10^{2}$. Abaixo
de $50$ quando $1.2\times 10^{5}e^{-0.5t} = 50$: $t = 2\ln(2400) \approx
\qty{16}{d}$.
\end{solution}

\begin{solution}{exo:b3:virology:7}
$uL = 0.36$; fração sem erro $e^{-0.36} = 0.70$. Com o triplo:
$uL =
1.08$, $e^{-1.08} = 0.34$ --- a maioria da prole passa a carregar
mutações e a sequência mais apta é diluída rumo ao limiar. Um
coronavírus de \qty{30}{kb} com $u = 10^{-4}$ teria $uL = 3$ e apenas
\qty{5}{\%} de cópias sem erro, além do limiar; sua exonuclease leva $u$
a $10^{-5}$, $uL = 0.3$, e três quartos ficam sem erro.
\end{solution}

\begin{solution}{exo:b3:virology:8}
Deriva: as mutações pontuais na hemaglutinina e na neuraminidase se
acumulam sob a seleção dos anticorpos, de modo que a linhagem de cada
temporada é em parte nova e a imunidade do ano passado fica em parte
obsoleta. Salto: a coinfecção de uma mesma célula por uma linhagem
humana e uma animal rearranja os oito segmentos, e um subtipo de
hemaglutinina de aves ou de porcos aparece num \omterm{def:b3:virology:virus}{vírus} adaptado aos
humanos. Ninguém tem anticorpos contra o subtipo novo, de modo que a
população inteira fica suscetível de uma vez --- a condição para uma
pandemia, cumprida em 1918, 1957, 1968 e 2009, a cada vez com uma
hemaglutinina nova.
\end{solution}

\begin{solution}{exo:b3:virology:9}
A timidina-cinase do herpesvírus fosforila o aciclovir, no qual as
cinases celulares mal tocam; as cinases celulares então completam o
trifosfato, que a DNA-polimerase viral incorpora e que, sem hidroxila
3$'$, termina a cadeia. As células não infectadas nunca fazem a forma
ativa, e a polimerase viral a prefere à celular. Resistência: perda ou
alteração da timidina-cinase viral (a mais comum), ou mutações na
polimerase viral que excluam o análogo.
\end{solution}

\begin{solution}{exo:b3:virology:10}
Com $T$ constante, as equações são lineares em $(I, V)$ com matriz
$\begin{pmatrix} -\delta & \beta T \\ p & -c \end{pmatrix}$, traço
$-(\delta + c) < 0$ e determinante $\delta c - \beta Tp$. Os autovalores
têm traço negativo, de modo que um deles é positivo exatamente quando o
determinante é negativo: $\beta Tp > c\delta$, isto é,
$R_{0} = \beta
Tp/(c\delta) > 1$. Fatores: $p/\delta$ é o número de
\omterm{def:b3:virology:virus}{vírions} que uma célula infectada faz em sua vida; $\beta T/c$ é o número
de células que um \omterm{def:b3:virology:virus}{vírion} infecta em sua vida; o produto é o número de
células infectadas que uma célula infectada produz.
\end{solution}

\begin{solution}{exo:b3:virology:11}
$10^{6} = 2^{20}$: vinte meias-vidas, $880$ meses, cerca de $73$ anos
--- mais que uma vida. Os fármacos que bloqueiam a replicação não tocam
num provírus que não está sendo transcrito; o reservatório persiste e,
quando a terapia para, uma única célula que reative reinicia a
infecção. Uma cura precisa remover ou silenciar permanentemente o
reservatório: matar as células latentes (reativá-las sob terapia para
que morram ou sejam mortas), excisar ou inutilizar o provírus, ou
substituir o sistema imune por um que o \omterm{def:b3:virology:virus}{vírus} não consiga infectar.
\end{solution}

\begin{solution}{exo:b3:virology:12}
$uL = 0.03$ mutação por genoma; $10^{9}$ genomas por dia carregam
$3\times
10^{7}$ mutações entre $9\times 10^{4}$ mutações únicas
possíveis: cada uma surge cerca de $300$ vezes por dia. Um mutante
duplo específico surge a $10^{-12}$ por genoma, $10^{-3}$ por dia ---
uma vez a cada três anos; um triplo a $10^{-18}$, nunca. Dois fármacos
com mutações de resistência independentes bastariam em princípio, e
três dão margem para adesão imperfeita. Dois fármacos que se ligam à
mesma enzima podem ser derrotados por uma única mutação que altere o
sítio ou a conformação da enzima, e suas mutações de resistência não
são independentes; eles contam como menos de dois.
\end{solution}

\begin{solution}{pb:b3:virology:1}
\textbf{1.} $2\times 10^{5}\times 1.5\times 10^{4} = 3\times 10^{9}$
\omterm{def:b3:virology:virus}{vírions}.
\textbf{2.} Depurados $cV = 3\times 3\times 10^{9} \approx 10^{10}$ por
dia; produzidos, os mesmos.
\textbf{3.} $I^{*} = cV^{*}/p = 9\times 10^{9}/500 = 1.8\times 10^{7}$
células infectadas; morrendo $\delta I^{*} = 9\times 10^{6}$ por dia.
\textbf{4.} RNA $2\times 9700\times 330 = 6.4\times 10^{6}$ Da; proteína
$2000\times 25\,000 = 5\times 10^{7}$ Da; total $5.6\times 10^{7}$ Da
$\approx 10^{-16}$ g. Todos os \omterm{def:b3:virology:virus}{vírions} juntos: $3\times 10^{9}\times
10^{-16} = 3\times 10^{-7}$ g, um terço de micrograma --- três mil
vezes menos que um grão de areia.
\textbf{5.} $1.8\times 10^{7}/10^{12} = 2\times 10^{-5}$ do conjunto.
Dez anos de $9\times 10^{6}$ mortes por dia dão $3\times 10^{10}$
células, três por cento do conjunto: o conjunto falha não por subtração
aritmética, mas porque uma década de regeneração forçada e de ativação
crônica esgota a capacidade de repô-las.
\textbf{6.} $9\times 10^{6}/24 \approx 4\times 10^{5}$ corpos sendo
removidos a cada momento.
\textbf{7.} $V = 2\times 10^{5}(3e^{-0.5t} - 0.5e^{-3t})/2.5$: $t = 1$:
$1.4\times 10^{5}$; $t = 3$: $5.4\times 10^{4}$; $t = 7$: $7\times
10^{3}$; $t = 14$: $2\times 10^{2}$.
\textbf{8.} $2.4\times 10^{5}e^{-0.5t} = 50$: $t = 2\ln(4800) \approx 17$
dias.
\textbf{9.} Os dias 3--10 caem na fase lenta: a inclinação de $\ln V$ ali
é $-\delta$ (de $5.4\times 10^{4}$ a $1.6\times 10^{3}$ em sete dias,
$\delta = 0.50$). A fase rápida dura apenas horas ($1/c$), de modo que
$c$ exige amostras de poucas em poucas horas no primeiro dia, ajustadas
à fórmula de duas exponenciais.
\textbf{10.} As células infectadas latentes, que liberam \omterm{def:b3:virology:virus}{vírus} à medida
que reativam; sua taxa de decaimento é $\ln 2/44$ por mês, cerca de
$5\times 10^{-4}$ por dia.
\textbf{11.} \omterm{def:b3:virology:virus}{Vírion} $\ln 2/3 = \qty{0.23}{d}$, cerca de \qty{5.5}{h};
célula infectada $\ln 2/0.5 = \qty{1.4}{d}$.
\textbf{12.} Uma carga constante parecia um \omterm{def:b3:virology:virus}{vírus} que não fazia nada; as
equações a mostram como um estado estacionário em que $10^{10}$ \omterm{def:b3:virology:virus}{vírions}
são feitos e depurados e $10^{7}$ linfócitos~T infectados e mortos todo
dia.
\textbf{13.} $9700\times 10^{-5} \approx 0.1$ mutação por genoma;
$10^{10}\times 0.1 = 10^{9}$ mutações por dia.
\textbf{14.} $3\times 9700 \approx \num{29000}$ mutações únicas
possíveis; cada uma surge $10^{9}/\num{29000} \approx 3\times 10^{4}$
vezes por dia.
\textbf{15.} Duplo específico: $10^{-10}$ por genoma, cerca de um por
dia sem tratamento; triplo específico: $10^{-15}$, $10^{-5}$ por dia ---
uma vez a cada três séculos.
\textbf{16.} $10^{7}\times 10^{-15} = 10^{-8}$ por dia,
$4\times 10^{-6}$ por ano.
\textbf{17.} Sem um dos fármacos, bastam mutantes resistentes aos
outros dois: um duplo específico a $10^{-10}$; $7\times 10^{7}$ genomas
na semana dão $7\times 10^{-3}$ esperados --- ainda improvável, a menos
que a carga suba de volta rumo a $10^{10}$ por dia durante o lapso,
quando isso passa a ser um por dia e o escape fica provável. As doses
esquecidas são o modo como a resistência chega.
\textbf{18.} Os análogos compartilham o sítio ativo da transcriptase
reversa, e uma única mutação ali (ou um conjunto de mutações de análogo
de timidina) altera o modo como a enzima lida com vários deles: a
resistência cruzada. Dois nucleosídeos contam como menos de dois
fármacos independentes, e por isso os esquemas combinam classes --- um
nucleosídeo, um inibidor de integrase, um \omterm{def:b3:virology:drugs}{inibidor de protease} ou um não
nucleosídeo.
\textbf{19.} $20$ meias-vidas, $880$ meses, $73$ anos; até $10^{4}$,
$\log_{2}100 = 6.6$ meias-vidas, $290$ meses, $24$ anos.
\textbf{20.} De um \omterm{def:b3:virology:virus}{vírion} a $3\times 10^{9}$ com $R_{0} = 8$ por
geração: $\ln(3\times 10^{9})/\ln 8 = 10.5$ gerações, três semanas.
\textbf{21.} Choque e morte: reativar as células latentes sob terapia
para que morram ou sejam mortas --- obstáculo: nenhum agente reativa
todas elas, e as células reativadas não são mortas de modo confiável.
Edição gênica: excisar o provírus ou destruir o receptor CCR5 nas
células do paciente --- obstáculo: chegar a cada célula. (Substituir a
medula por células de doador sem CCR5 curou um punhado de pacientes, a
um risco aceitável apenas quando o transplante já era necessário de
qualquer modo.)
\textbf{22.} $1 - 1/R_{0}$: \qty{67}{\%} para $R_{0} = 3$ e
\qty{80}{\%} para $5$.
\textbf{23.} A terapia faz $\beta \to 0$ dentro do paciente, a carga cai
pelo teorema a quase nada, e a transmissão, que é proporcional à carga,
cessa; tratar cada pessoa infectada leva o $R_{0}$ da população abaixo
de um.
\textbf{24.} $f = (1 - 1/R_{0})/E$: com $R_{0} = 4$ e $E = 0.7$,
$f =
0.75/0.7 = 1.07$ --- mais que todo mundo: a \omterm{def:b3:virology:drugs}{vacina} sozinha não
consegue alcançar o limiar; com $E = 0.9$, $f = 0.83$.
\textbf{25.} Cerca de $10^{10}$ \omterm{def:b3:virology:virus}{vírions} produzidos por dia;
indetectável por volta do dia $17$; uns $4\times 10^{-6}$ mutantes
triplos por ano em terapia.
\end{solution}
